% Kapitel Diskussion der Hausarbeit
% ----------------------------------------------------------------------
% GERUEST-Konventionen: wie sections/theorie.tex (SLOT/NOTIZ/TODO).
% Vor Abgabe muessen alle SLOT/NOTIZ/TODO-Marker aufgeloest sein.

% SLOT D1 — Einordnung der Befunde (2-3 eigene Saetze)
% Behauptung: den H1-Nullbefund im modernen Zeitraum (1992-2023) vor dem
% Hintergrund des Originals einordnen (dort ca. 30 Prozent weniger
% Konditionen). Moegliche Erklaerungen, abwaegen: veraenderte
% Konditionalitaetspraxis nach 2008, zu kleine UNSC-Zellen (Power-Problem),
% oder tatsaechlich nachlassende Geberlogik der Hauptanteilseigner.
% WICHTIG: die Fuel-Export-Ausnahme (p = 0,024) nicht ueberbewerten —
% eine einzige Spezifikation von vielen (Mineral-Variante p = 0,91,
% Placebo gesamte Exporte p = 0,37). Leave-one-out existiert NUR fuer
% die HAUPTinteraktion resource_dep (dort: Abhaengigkeit von PER/GAB);
% fuer die Fuel-Spezifikation selbst liegt kein LOO-Befund vor.
% NOTIZ (eigenes Phase-A-Ergebnis — staerkstes Argument hier):
% Auf den ORIGINALDATEN gilt (s. docs/session_protokoll_2026-09-24.md
% und Ergebnistabelle original_replication_zeitfenster; Zahlen vor
% Abgabe gegen die Tabelle pruefen):
%   1992-2001: unsc3 = -4,65 (p = 0,042) — Effekt signifikant.
%   2002-2008: unsc3 = -1,86 (p = 0,52) — statistisch Null.
% D.h. der publizierte Effekt wird von den 1990ern getragen und ist
% ab 2002 selbst auf Originaldaten null. Der Phase-B-Nullbefund
% (1992-2023) ist damit KEIN Replikationsfehler, sondern bildet einen
% realen zeitlichen Rueckgang ab. Achtung fuer die Formulierung: Das
% Originalpapier selbst berichtet KEINE Subperioden-Zerlegung — die
% Einsicht stammt aus der eigenen Phase-A-Analyse und muss so
% (als eigener Befund) zitiert werden, nicht DSV zugeschrieben werden.
% ERGAENZUNG (direkter Differenztest, phase0_replikation_original.R):
% Interaktion unsc3 x ab2002 — Basis: +2,54 (SE 2,74, p = 0,36);
% Vollmodell: +5,21 (SE 5,42, p = 0,34); Tabelle
% results/phase_a/tables/original_replication_zeitfenster_interaktion.csv.
% Richtung wie erwartet (Effekt nach 2002 deutlich schwächer Richtung
% Null), aber die Differenz ist auf dem kleinen Sample NICHT
% signifikant. In der Arbeit ehrlich beides berichten: signifikante
% Fruehperiode + Null spät + nicht-signifikanter Differenztest.
% NOTIZ (Finanzkrise/IMF-Rueckkehr, verifiziert gegen Kentikelenis et al.
% 2016 — alle drei Stellen auf ARTIKEL-S. 545, Tab.-1-Diskussion ca. S. 551):
% - Mitte der 2000er: Nachfrage nach IMF-Krediten auf historischem Tief,
%   Zweifel an der Zukunft des IMF (dort zitieren die Autoren Momani und
%   Helleiner 2007 — Original NICHT in literatur.bib: entweder Bib-Eintrag
%   + eigene Lektuere nachtragen ODER als Sekundaerzitat nur ueber
%   Kentikelenis et al. laufen lassen).
% - 2009: G-20-Zusagen von 750 Mrd. USD; 2009-2014 insgesamt 129 neue
%   Kredite an 76 Laender (Artikel-S. 545).
% - Post-2008: Konditionen-Median sinkt kurzfristig (von ca. 42 in
%   1996-2007 auf 33-35), steigt bis 2014 wieder auf 44; Zahl aktiver
%   Programme nimmt nach der Krisenspitze wieder ab.
% Verwendung im Argument: Der Phase-B-Nullbefund ist KEIN Artefakt der
% IMF-Inaktivitaet — auch in der Phase massiver Nachkrisen-Kreditvergabe
% zeigt sich kein UNSC-Effekt im erweiterten Panel. Caveat: Nachkrisen-
% programme unterscheiden sich in Typ und Zusammensetzung von den
% BOP-Krisen der 1990er; die Einflusskauf-Logik war an deren geopolitische
% Konstellation geknuepft (vgl. Originalpapier: Zimbabwe/Irak 1992).
% Quellen: \parencite[TODO]{DreherPoliticsIMFConditionality2015},
% \parencite[TODO]{KentikelenisIMFConditionalityDevelopment2016},
% \parencite[TODO]{StubbsHowEvaluateEffects2020} (Endogenitaet,
% Identifikation, Evaluationsproblematik).

% SLOT D2 — Limitationen buendeln (2 eigene Saetze)
% Behauptung: die Daten-/Klassifikationslimitationen (SLOT M1/M2 in
% daten-methodik.tex) hier als Methodik-Diskussion zusammenfuehren, nicht
% als reine Fehlerliste: was koennen die Modelle deshalb aussagen und was
% nicht (deskriptiv-vorsichtig, keine kausale Lesart).

% SLOT D3 — Abhaengigkeitstheoretische Lesart (1-2 eigene Saetze)
% Behauptung: aus der Rahmung in theorie.tex (SLOT 4/5) folgt: ein
% ausbleibendes "Entgegenkommen" fuer rohstoffreiche Laender (H2-Null,
% H4-Null) spricht gegen eine Machtressourcen-Logik im modernen
% Konditionalitaetsregime — oder dafuer, dass deren Mechanismus nach 2008
% an Bedeutung verloren hat. Formulierung braucht Sorgfalt, weil H2/H4
% nur explorativ getestet wurden.
% Quellen: \parencite[TODO]{MuchhalaExaminingFeaturesEconomic2025},
% \parencite[TODO]{AntunesDeOliveiraDependencyTheory2024}.
% NOTIZ (fakultativ, 1 Satz): Auch auf der Outcome-Seite kein Muster —
% Konditionalitaetsintensitaet korreliert nicht mit der Veraenderung
% der Rohstoffexportabhaengigkeit t bis t+3 (Spearman rho = 0,039,
% n = 333 Programm-Landjahre; h4_conditionality_resource_change_3y):
% kein Beleg dafuer, dass Konditionalitaet Rohstofforientierung
% systematisch auf- oder abbaut — in keine Richtung.

\subsection{Beitrag der Arbeit}
% SLOT D4 — Beitrag (2-3 eigene Saetze)
% Behauptung: die Replikation uetraegt das DSV-Design auf den modernen
% Zeitraum 1992-2023 und zeigt: der Originalbefund laesst sich nicht
% ohne weiteres fortsetzen. Explorative H3/H4-Ergebnisse liefern
% Hypothesen fuer Folgestudien (Kompositionshypothese konfirmatorisch
% pruefen, sobald eine belastbarere Policy-Area-Klassifikation vorliegt).

% SLOT D5 — Ausblick (1-2 eigene Saetze, optional)
% Behauptung: konfirmatorische Erweiterungen nennen: eigene
% Policy-Area-Klassifikation fuer post-2008-Bedingungen, groessere
% Zahl UNSC-Fenster im Panel, alternativ IMF-Kreditbestand statt
% Nettofluesse (bereits gerechnet, s. Robustheitstabelle).
